% Draft of the numerical section, written against the macros already defined in
% the manuscript preamble (\Id, \N, \Aspace, \Delta t, \Fobj, \post, \grad, \norm,
% \E) and the existing labels \eqref{upd:Riem}, \eqref{upd:KL}, \eqref{ODEs:NG},
% Proposition~\ref{prop:affine-equivariance}, Theorem~\ref{thm:Riemannian-global},
% Theorem~\ref{thm:KL-global}, Proposition~\ref{prop:local-linearized-generator},
% Theorem~\ref{thm:local-asymptotic-rates}, Theorem~\ref{thm:global-to-local}.
%
% Every number quoted below is read from results/processed/manuscript/*.csv and
% results/tables/manuscript_*.csv.  Figures are results/figures/manuscript/figure_{1,2,3}.pdf.

\section{Numerical experiments}
\label{sec:numerics}

The experiments in this section illustrate the mechanisms behind the main
theorems rather than benchmark an implementation. They are organized as three
groups: the geometric structure of the Gaussian problem, global-to-local
convergence on non-Gaussian targets, and a genuine variational-inference problem.

The study is deterministic throughout. The stochastic results of
Section~\ref{sec:stochastic} are non-asymptotic guarantees rather than an
empirical contribution, and are not studied systematically here.

\subsection{Experimental setup}
\label{subsec:numerics-setup}

\paragraph{Methods.}
Only three schemes iterate: the Riemannian retraction \eqref{upd:Riem}
(\emph{FR-R}), the KL/Bregman scheme \eqref{upd:KL} (\emph{FR-KL}), and the
deterministic forward--backward Bures--Wasserstein method of
\cite{diao2023forward}, written \emph{BW-FB} so that all three labels name their
geometry before their discretization; it is used only as an external baseline. On the logistic-regression problem the Laplace approximation is
reported as a non-iterative reference. No Wasserstein warm start, hybrid
geometry, or covariance rescue enters any deterministic comparison.

\paragraph{Expectations.}
Every expectation is computed exactly; no Monte Carlo enters the deterministic
study. Gaussian targets are closed form. The separable targets of
Section~\ref{subsec:numerics-nongaussian} use a panelled Gauss--Legendre rule
against the Gaussian density, which resolves the unit-width
$\operatorname{sech}^2$ structure uniformly in the marginal width. The logistic
posterior admits the same reduction: $V$ depends on $\theta$ only through the
linear predictors $z_i=x_i^{\!\top}\theta$, and $z_i\sim\N(x_i^{\!\top}m,
x_i^{\!\top}Cx_i)$ under $q$, so the objective, the expected gradient and the
expected Hessian are $n$ independent one-dimensional integrals. The updates use
order $48$ and the objective evaluation and reference solve use order $96$; the
two agree to about $10^{-11}$, so the evaluation rule is both independent of and
strictly finer than the update rule.

\paragraph{Reference solutions.}
A residual is a gradient norm and a gap is an energy difference, so the two are
related through the Gaussian variational PL inequality of
Lemma~\ref{lem:gaussian-PL}, $\Delta(a)\le\|g(a)\|^2+\operatorname{Tr}((C^{-1}-H)C(C^{-1}-H))$
divided by $2\alpha_\star$. This is a proved bound with a constant computed
exactly for each target family. Applied to each reference's own residual it
certifies the reference's suboptimality at least $7.4$ orders of magnitude below
the smallest gap any panel draws. The same two lemmas give a reference-free
certificate at every iterate,
$\Delta(a)\le\min\{\|\cdot\|_{\mathrm{BW}}^2/(2\alpha_\star),
\|\grad\mathcal E(a)\|_a^2/(\alpha_\star\lambda_{\min}(C))\}$, which holds at
every state where the measured gap is meaningful and is tight to a factor of
$1.2$ on the log-cosh grid.

\paragraph{Stepsizes.}
Panels that verify a theorem use the stepsize that theorem certifies. The
remaining panels use one frozen practical multiplier per method, selected once on
two pilot cells that appear in no figure and never revisited. The multiplier
applies to $1/(\beta_\star\max\{\lambda_{0,\star}^{\max},1\})$; a multiple of the
certified stepsize is not transferable, because the certified constant carries a
worst-case covariance-growth allowance that is not attained. The frozen values
are $2$ for both Fisher--Rao schemes and $1$ for BW-FB; on the optimizer-whitened
targets $\beta=\beta_\star$, so that multiplier is exactly the ceiling
\cite{diao2023forward} admits. Certified stepsizes are
reported alongside the stepsizes actually used.

\paragraph{Implementation.}
All computations are in double precision with NumPy and SciPy; SPD matrix
functions go through symmetric eigendecompositions and linear solves through
Cholesky factors. The campaign is $131$ trajectories produced from a single
tagged revision on a clean tree, run on one x86-64 node, and completes in minutes.

\paragraph{Reported quantities.}
We report the normalized objective gap $\Delta(a_n)/\Delta(a_0)$, the Fisher--Rao
gradient norm $\norm{\grad\mathcal E(a_n)}_{a_n}$, and the optimizer-whitened
distance $\norm{a_n-a_\star}_\star$, against iteration count, elapsed flow time
$n\Delta t$, and wall-clock time. Each deterministic iteration consumes one
expected gradient and one expected Hessian, so the iteration count is already an
oracle count. Curves are truncated where they reach the resolution of the
quantity being plotted, so no numerical plateau is drawn as convergence.

\subsection{Gaussian structure and affine invariance}
\label{subsec:numerics-gaussian}

Figure~\ref{fig:gaussian} collects the three measurements that motivate the
geometry.

Panel~(a) isolates the initialization term of
Theorem~\ref{thm:Riemannian-global}. The target is $\N(0,\Id)$ in $N_\theta=20$,
so $C_\star=\Id$, $\beta_\star=1$, and the covariance band
$C_n\succeq(2\beta_\star)^{-1}C_\star$ reduces to
$\lambda_{\min}(C_n)\geq\frac12$. Starting from $C_0=\lambda_0\Id$ with
$\lambda_0$ ranging over four decades and a common step inside both certified
windows, the entry time grows with unit slope in
$\log_+[1/(\beta_\star\lambda_{0,\star})]$: the retraction sits on the predicted
line with fitted slope $0.999$ and intercept $-0.000$, and the Bregman scheme
runs about five per cent above it. Over the four initializations measured this is
consistent with the predicted logarithmic dependence on the initial covariance.

Panel~(b) tests Proposition~\ref{prop:affine-equivariance} directly. A base
problem in $N_\theta=10$ is pushed through $\theta\mapsto S\theta+b$ with
$\operatorname{cond}(S)$ ranging over six decades, and each transformed
trajectory is mapped back to the base coordinates. The Fisher--Rao iterates
reproduce the base trajectory to roundoff, with a discrepancy that grows like
$\varepsilon_{\mathrm{mach}}\operatorname{cond}(S)^2$: from $10^{-15}$ at
$\operatorname{cond}(S)=1$ through $1.1\times10^{-12}$, $4.3\times10^{-9}$ and $5.9\times10^{-5}$. This is
floating-point amplification through an ill-conditioned map, not a loss of
equivariance; the exponent is $2$ because the covariance is carried by the
congruence $C\mapsto SCS^{\!\top}$. The grid stops at
$\operatorname{cond}(S)=10^{6}$: one decade further the transported covariance
has condition number $\varepsilon_{\mathrm{mach}}^{-1}$ and no method can be
measured. BW-FB, whose Bures--Wasserstein geometry is not affine invariant,
departs by an $O(1)$ amount, about $0.5$, as soon as $S$ is not orthogonal.

Panel~(c) shows the algorithmic consequence on the anisotropic Gaussian
benchmark of \cite{diao2023forward} in $N_\theta=10$ with logarithmically spaced
precision eigenvalues, where $\kappa=10^{9}$ but $\kappa_\star=1$. Each method
runs at its own certified stepsize. The Riemannian retraction reaches a
normalized KL gap of $10^{-12}$ after $58$ iterations and the Bregman scheme after
$47$; BW-FB reduces the gap only to $0.45$ of its initial value in $500$. The target is constructed so that coordinate
anisotropy and intrinsic conditioning differ maximally, and the panel should be
read as separating those two notions rather than as evidence that Fisher--Rao
dominates in general.

\subsection{Global and local behavior on non-Gaussian targets}
\label{subsec:numerics-nongaussian}

The target family is
\[
    V(\theta)
    =
    \tfrac12\theta^{\!\top}Q\theta
    +
    \rho\sum_j\log\cosh(b_j^{\!\top}\theta-c_j),
\]
which is strongly logconcave and smooth for every $\rho\geq0$. The affine map is
chosen so that $C_\star=\Id$ to numerical tolerance, which is possible because the
separable optimizer does not depend on it; the construction is verified to
$10^{-12}$ in Frobenius norm on every cell. Starting from $m_0=2e_1$, $C_0=\Id$ then places
the covariance inside the band from the first iteration, so the trajectory
measures non-Gaussian localization rather than a covariance burn-in. The grid is
$N_\theta\in\{10,50\}$, $\kappa_{\mathrm{base}}\in\{1,10,10^2\}$,
$\rho\in\{0.1,1\}$, one instance per cell.

Panels~(a)--(c) of Figure~\ref{fig:nongaussian} show the normalized gap on three
preregistered cells. All three methods contract at comparable per-iteration
rates; the two Fisher--Rao schemes are within a factor of two of one another
throughout, and the retraction is the faster of the two on the harder cells.

Panels~(d) and~(e) summarize the whole grid, and together they separate a
stepsize effect from a geometric one. In iteration count the methods differ by up
to an order of magnitude; in elapsed flow time they agree to within a factor of
$2.2$ on the eight cells with $\kappa_\star\geq2$. Two observations follow. First,
the time to tolerance does not grow with $\kappa_\star$ over two decades of it, so
the worst-case constant of Theorem~\ref{thm:global-to-local} is far from attained
on this family; the intrinsic conditioning enters the certificate but not the
realized rate. Second, on the four least-conditioned cells the single frozen
multiplier places $\Delta t\,\gamma_\star$ near $2$ for the Fisher--Rao schemes,
the linearized one-step factor $1-\Delta t\,\gamma_\star$ approaches $-1$, and the
iteration count rises by an order of magnitude. That is the cost of freezing one
number across a grid on which $\beta_\star$ varies by a factor of five, and not a
property of the flow.

Panel~(f) verifies Theorem~\ref{thm:local-asymptotic-rates} quantitatively. For
two targets of the same family we assemble the linearized generator
$\mathcal L_\star$ of Proposition~\ref{prop:local-linearized-generator}, take the
eigenvector $v_{\min}$ of its smallest eigenvalue $\gamma_\star$, and initialize
at $a_0=a_\star+rv_{\min}$ for $r\in\{10^{-1},5\times10^{-2},10^{-2}\}$. The
measured per-iteration contraction is compared with $1-\Delta t\,\gamma_\star$ for
\eqref{upd:Riem} and with $1-\Delta t\,\gamma^{\mathrm{KL}}_{\star,\Delta t}$ for
\eqref{upd:KL}. The three radii agree with one another to five significant
figures, so the measured rate is a property of the mode and not of the initial
amplitude, and the residual discrepancy against the prediction is the
second-order term of the discretization: at most $3.5\times10^{-5}$ on the
$\rho=1$ targets and $1.1\times10^{-3}$ on the $\rho=0.1$ targets, whose certified
step is five times larger.

\subsection{Bayesian logistic regression}
\label{subsec:numerics-logistic}

The last group is a genuine variational-inference problem: synthetic Bayesian
logistic regression with a proper Gaussian prior
$\theta\sim\N(0,\lambda^{-1}\Id)$, $\lambda=1$, in $N_\theta=50$ with $500$
training and $5000$ held-out points, at feature conditionings
$\kappa_X\in\{1,10^2,10^4\}$ and five independently generated datasets each, one
master seed per dataset. All methods start from the prior, $m_0=0$,
$C_0=\lambda^{-1}\Id$, which overshoots $C_\star$ by about two orders of magnitude
and so exercises the covariance burn-in that the whitened targets of
Section~\ref{subsec:numerics-nongaussian} deliberately avoid.

Panels~(a) and~(b) of Figure~\ref{fig:logistic} show the objective gap at
$\kappa_X=10^4$ against iteration and against the time spent inside the update.
All three schemes converge geometrically over more than fourteen decades. The
terminal Fisher--Rao gradient norms are below $10^{-11}$ at $\kappa_X=10^4$ and
below $3\times10^{-8}$ at $\kappa_X=10^2$; at $\kappa_X=1$ the Fisher--Rao schemes
are still at $2.5\times10^{-3}$ after $2000$ iterations, which is where their
stepsize is smallest.

No method dominates. The median iteration count to a $10^{-6}$ relative gap is
\[
\begin{array}{lccc}
 & \kappa_X=1 & \kappa_X=10^2 & \kappa_X=10^4\\
\text{FR-R} & 1180 & 490 & 360\\
\text{FR-KL} & 1180 & 500 & 370\\
\text{BW-FB} & 60 & 360 & 470.
\end{array}
\]
BW-FB is an order of magnitude faster on the best-conditioned problem and slower
on the worst, and the crossover is near $\kappa_X=10^2$. This is the expected
shape: the Bures--Wasserstein guarantee places the conditioning in the rate and
the Fisher--Rao guarantee places it in the admissible stepsize, so the two trade
places as the conditioning grows. The two Fisher--Rao schemes agree with one
another to within three per cent throughout.

The cost per iteration is the same for all three methods to within eight per
cent, $12.3$ to $14.1$ milliseconds, so panel~(b) repeats the shape of panel~(a).
The cost is dominated by the expectation, which is $O(nQ)$ for $n=500$
observations and $Q$ quadrature nodes, rather than by the $O(N_\theta^3)$ linear
algebra that separates the matrix exponential of \eqref{upd:Riem} from the
resolvent solve of \eqref{upd:KL}. That distinction would become visible in
wall-clock terms only at large $N_\theta$ with a cheaper expectation.

The held-out predictive negative log-likelihood at termination is reported in
Table~\ref{tab:logistic-predictive}. All three iterative schemes agree to four
decimals on every conditioning and every dataset, as they must, since they
converge to the same optimizer; the differences between them are in how fast they
reach that common answer, not in the answer. The Laplace approximation, which is
a fixed point of none of the schemes, attains a \emph{lower} predictive negative
log-likelihood than the Gaussian variational optimum at all three conditionings,
by $0.0083$, $0.0047$ and $0.0012$ respectively, while its objective gap is larger
by many orders of magnitude. Minimizing the reverse Kullback--Leibler divergence
over Gaussians is therefore not the same as optimizing held-out predictive
quality on this problem; we report the observation without attributing it to a
mechanism, which these experiments do not isolate.
